\section{Experiments \& Results}
\label{sec:exp}

\textbf{Partitioning.} Unless stated otherwise, every experiment uses the
five predefined subject-wise folds (60/20/20 subjects). Training uses the 60
training subjects, hyper-parameters and the decision threshold are selected on
the 20 validation subjects, and the 20 test subjects are predicted once. The
five test sets are disjoint, so pooling them yields exactly one prediction per
subject and 10\,407 window predictions in total. This is cross-validation in
the sense that every subject is tested exactly once, without any subject
appearing in two roles.

\textbf{Experiments performed.} (1) token representation: raw samples versus
amplitude statistics versus stride cycles versus their fusion; (2) architecture:
unidirectional versus bidirectional, three encoder sizes, regularisation and
seed ensembling; (3) foot configuration: left, right, both; (4) protocol
comparison: subject-wise versus window-wise splits. Fifteen GPU runs were
executed in total.

\subsection{Experiment 1: token representation}
Table~\ref{tab:tokens} reports subject-level accuracy on both feet for the
input representations. Raw signals were markedly unstable: on folds 1 and 4
the validation AUC stayed near chance ($0.53$ and $0.54$) while folds 2, 3 and
5 exceeded $0.92$, and neither longer windows (30\,s, 60\,s), seed restarts
nor a larger encoder repaired it. Amplitude-statistic tokens removed the
instability -- the weakest fold rose to $0.75$ validation AUC -- and improved
accuracy by seven points. Stride-cycle tokens alone were weaker, and fusing
them with amplitude tokens did not surpass amplitude tokens alone, although at
subject level the two families were complementary in a Random Forest analysis
(AUC $0.79$ stride only, $0.83$ amplitude only, $0.86$ combined).

\begin{table}[t]
  \centering
  \caption{Effect of the input representation on subject-level accuracy (both
  feet, subject-wise protocol). ``Worst fold'' is the lowest validation ROC-AUC
  over the five folds and exposes the training collapse seen with raw input.}
  \label{tab:tokens}
  \begin{tabular}{lcc}
    \toprule
    Input representation & Accuracy (\%) & Worst fold AUC \\
    \midrule
    Raw samples, 5\,s windows            & 69 & 0.53 \\
    Raw samples, 30\,s windows           & 69 & 0.63 \\
    Raw samples, 60\,s windows           & 65 & 0.59 \\
    Stride-cycle tokens                  & 66 & 0.71 \\
    Amplitude ++ stride (fused)          & 78 & 0.71 \\
    \textbf{Amplitude statistics}        & \textbf{80} & \textbf{0.77} \\
    \bottomrule
  \end{tabular}
\end{table}

\subsection{Experiment 2: architecture ablation}
With amplitude tokens fixed, encoder size was the decisive architectural
factor: enlarging the encoder from $d_{\text{model}}=128$, $L=4$ to
$d_{\text{model}}=256$, $L=6$ raised subject-level accuracy from 72\,\% to
80\,\%.

The bidirectional extension did \emph{not} reproduce that gain. At
$d_{\text{model}}=128$ the causal and the bidirectional encoder reach the same
calibrated accuracy (72\,\% each); the causal variant is in fact slightly
better at the fixed threshold (79\,\% versus 72\,\%) and in ranking quality
(AUC $0.87$ versus $0.82$). Since a bidirectional scan doubles the mixer
parameters, part of the size effect and the directionality effect are
confounded in the larger model, and we do not claim a benefit from
bidirectionality on this dataset. This is a negative result worth stating: the
intuition that a classifier over a fully observed window should see both
directions is not supported at the sample size available here.

Models overfitted quickly -- the best epoch was typically between 1 and 4 --
so the reported model uses dropout up to $0.3$, cosine decay with warm-up and
three-seed averaging; heavier regularisation (dropout $0.5$) did not help
further. Seed averaging lowers the peak (the best single seed reached
80\,\%) but is the more honest estimate, single seeds ranging from 72\,\% to
80\,\%.

\subsection{Experiment 3: comparison with the baselines}
Table~\ref{tab:main} is the main result: all models under the subject-wise
protocol, evaluated with identical code and aggregation. The Mamba model is
best on both feet (78\,\%) and on the right foot (79\,\%); on the left foot
the Random Forest remains slightly ahead (74\,\% versus 70\,\%). Sensitivity
is consistently above specificity for Mamba (82\,\% versus 74\,\% on both
feet), i.e.\ it misses few patients, which is the preferable error profile for
screening. Figure~\ref{fig:results}(a) visualises the comparison, and
Figure~\ref{fig:confusion} shows the confusion matrices.

\begin{table}[t]
  \centering
  \caption{Subject-level results under the subject-wise (leakage-free)
  protocol, pooled over the five disjoint test folds (100 subjects); mean
  $\pm$ standard deviation over 1000 bootstrap resamples. Precision, recall
  and F1 are class-support weighted, so recall equals accuracy by
  construction. PD is the positive class.}
  \label{tab:main}
  \setlength{\tabcolsep}{3.2pt}
  \begin{tabular}{llccccc}
    \toprule
    Model & Data & Acc (\%) & Pre (\%) & F1 & Sen (\%) & AUC \\
    \midrule
    \multirow{1}{*}{SVM-RBF} & Left  & 64 $\pm$ 5 & 65 $\pm$ 5 & 0.64 & 60 & 0.74 \\
                             & Right & 56 $\pm$ 5 & 57 $\pm$ 5 & 0.56 & 54 & 0.65 \\
                             & Both  & 68 $\pm$ 5 & 69 $\pm$ 5 & 0.68 & 66 & 0.77 \\
    \midrule
    \multirow{1}{*}{1D-ResNet} & Left  & 68 $\pm$ 5 & 69 $\pm$ 5 & 0.68 & 62 & 0.71 \\
                               & Right & 61 $\pm$ 5 & 65 $\pm$ 5 & 0.59 & 38 & 0.70 \\
                               & Both  & 66 $\pm$ 5 & 67 $\pm$ 5 & 0.66 & 56 & 0.76 \\
    \midrule
    \multirow{1}{*}{LSTM} & Left  & 70 $\pm$ 4 & 71 $\pm$ 4 & 0.70 & 72 & 0.73 \\
                          & Right & 68 $\pm$ 5 & 68 $\pm$ 5 & 0.68 & 68 & 0.69 \\
                          & Both  & 63 $\pm$ 5 & 64 $\pm$ 5 & 0.63 & 54 & 0.68 \\
    \midrule
    \multirow{1}{*}{Random Forest} & Left  & \textbf{74 $\pm$ 4} & 75 $\pm$ 4 & \textbf{0.74} & 76 & \textbf{0.83} \\
                                   & Right & 78 $\pm$ 4 & 79 $\pm$ 4 & 0.78 & 75 & 0.81 \\
                                   & Both  & 73 $\pm$ 4 & 74 $\pm$ 4 & 0.73 & 74 & 0.82 \\
    \midrule
    \multirow{1}{*}{\textbf{Mamba (ours)}} & Left  & 70 $\pm$ 4 & 71 $\pm$ 4 & 0.70 & 76 & 0.79 \\
                                  & Right & \textbf{79 $\pm$ 4} & \textbf{80 $\pm$ 4} & \textbf{0.79} & 84 & 0.82 \\
                                  & Both  & \textbf{78 $\pm$ 4} & \textbf{79 $\pm$ 4} & \textbf{0.78} & 82 & 0.82 \\
    \bottomrule
  \end{tabular}
\end{table}

\begin{figure}[t]
  \centering
  \includegraphics[width=\columnwidth]{figures/fig3_results}
  \caption{(a) Subject-level accuracy of all models under the subject-wise
  protocol, per foot configuration. (b) The same models on both feet under the
  two protocols: replacing subject-wise splits by window-wise splits inflates
  accuracy by roughly twenty points without any change to features, models or
  metrics.}
  \label{fig:results}
\end{figure}

\begin{figure}[t]
  \centering
  \includegraphics[width=\columnwidth]{figures/fig4_confusion}
  \caption{Subject-level confusion matrices of the final Mamba model (100
  subjects, pooled over the five disjoint test folds), for the left foot, the
  right foot and both feet. Rows are true classes, columns predicted classes;
  HC = healthy control, PD = Parkinson's disease.}
  \label{fig:confusion}
\end{figure}

\subsection{Experiment 4: the evaluation protocol}
\label{ssec:protocols}
To quantify the influence of the split unit, we repeated the evaluation with
windows -- rather than subjects -- assigned at random to the 60/20/20 splits,
keeping the data, features, models, metrics and code identical
(Figure~\ref{fig:protocols}). Table~\ref{tab:protocols} contrasts the two.
Accuracy rises by twenty points or more, and the Random Forest reaches
98\,\%. Under this protocol a per-subject decision is meaningless, because a
subject's windows are distributed across training and test, so only
window-level numbers can be quoted. \emph{These numbers measure how well a
model recognises previously seen individuals, not whether it detects
Parkinson's disease in a new patient}, and we report them solely as a
reference point for the literature.

\begin{figure}[t]
  \centering
  \includegraphics[width=\columnwidth]{figures/fig2_protocols}
  \caption{The two evaluation protocols. Each row is a subject, each cell one
  window. (a) Subject-wise: whole subjects are held out, so nothing about a
  test subject is seen in training. (b) Window-wise: test windows are drawn at
  random, so most test windows have highly correlated neighbours from the same
  recording in the training set.}
  \label{fig:protocols}
\end{figure}

\begin{table}[t]
  \centering
  \caption{Effect of the split unit on both feet, with everything else held
  fixed. Subject-wise values are subject-level; window-wise values are
  window-level, because per-subject decisions do not exist when a subject's
  windows are split across training and test. The right-hand column is
  reported for contrast only and must not be compared with the literature's
  subject-wise results.}
  \label{tab:protocols}
  \begin{tabular}{lcc}
    \toprule
    Model & Subject-wise (\%) & Window-wise, leaky (\%) \\
    \midrule
    SVM-RBF        & 68 & 88 \\
    Random Forest  & 73 & 98 \\
    Mamba (ours)   & \textbf{78} & 97 \\
    \bottomrule
  \end{tabular}
\end{table}
